\begin{enumerate}
  \item Les respostes s'han d'ajustar a l'enunciat de la pregunta. Es valora sobretot que l'alumnat demostri que té clars els conceptes de caràcter físic sobre els quals tracta cada pregunta.
  \item Es té en compte la claredat en l'exposició dels conceptes, dels processos, dels passos que cal seguir i de les hipòtesis, l'ordre lògic, l'ús correcte dels termes científics i la conceptualització segons l'enunciat.
  \item En les respostes cal que l'alumnat mostri una adequada capacitat de comprensió de les qüestions plantejades i organitzi la resposta de manera lògica, tot analitzant i utilitzant les variables en joc. També es valora el grau de pertinença de la resposta, el que l'alumnat diu i les mancances manifestes sobre el tema en qüestió.
  \item Les respostes s'han de raonar i justificar. Un resultat erroni amb un raonament correcte es valora. Una resposta correcta sense raonament ni justificació es pot valorar amb un 0.
  \item Un error no es penalitza dues vegades en el mateix problema. Si un apartat necessita un resultat anterior, i aquest resultat és erroni, es valora la resposta independentment del seu valor numèric i es té en compte el procediment de resolució.
  \item Si la resolució presentada a l'examen és diferent però correcta i s'ajusta al que es demana a l'enunciat, s'avalua positivament, encara que no coincideixi amb la resolució donada a la pauta de correcció.
  \item Un o més errors d'unitats o no posar-les (resultats intermedis i finals) en un problema es penalitzen amb 0,25 punts en aquest problema.
  \item Cal resoldre els exercicis fins al final i no es poden deixar indicades les operacions.
  \item Cal fer la substitució numèrica en les expressions que s'utilitzen per resoldre les preguntes.
  \item Un o més resultats amb un nombre molt elevat de xifres significatives (6 xifres significatives) o amb només una xifra significativa es penalitzen amb 0,10 punts en aquest problema.
\end{enumerate}
